% !TEX program = xelatex
\documentclass[12pt]{article}
\usepackage[margin=28mm]{geometry}
\usepackage{hyperref}
\usepackage{amsmath, amsthm, amssymb, amsfonts}
\usepackage{tikz-cd}
\usepackage{fontspec}
\usepackage{unicode-math}
\setmainfont{Latin Modern Roman}
\setmonofont{Latin Modern Mono}
\usepackage{listings}
\lstset{basicstyle=\ttfamily\footnotesize, columns=fullflexible, breaklines=true, showstringspaces=false}

\title{ClaudeD\_Kleisli — Lean→TeX Export (English)}
\author{TeX generated by ChatGPT-5 Pro\\Lean source generated by CODEX (OpenAI)}
\date{2025-09-09}

\begin{document}
\maketitle

\begin{abstract}
Minimal Kleisli normalization for `productMonad` (Δ ⊣ ×).
  - Stays entirely inside the Kleisli category; no mixing with TopCat arrows.
  - Uses definitional (`rfl`) lemmas so `ext; simp [Category.assoc]` cooperates
    with the η/μ component lemmas in ClaudeD.
\end{abstract}

\section*{Provenance}
This \LaTeX{} file was generated from the Lean source file \texttt{ClaudeD\_Kleisli.lean}.\\
\textbf{TeX generation}: ChatGPT-5 Pro.\\
\textbf{Lean source generation}: CODEX (OpenAI).

\section*{At-a-glance}
Total lines: 27\\
Defs: 0\\
Lemmas: 0\\
Theorems: 0\\
Structures: 0\\
Inductives: 0

\section*{Sample diagrams}

\[
\begin{tikzcd}
A \arrow[r, "f"] \arrow[d, "g"'] & B \arrow[d, "h"] \\
C \arrow[r, "k"'] & D
\end{tikzcd}
\]



\[
\begin{tikzcd}
A \arrow[r, "f: A \to TB"] \arrow[rr, bend left=15, "{g \star f}"] & TB \arrow[r, "{\mu_B \circ T g}"] & TC
\end{tikzcd}
\]


\section*{Declarations (first-line signatures)} 
\begin{center}
\begin{tabular}{lll}
\textbf{kind} & \textbf{name} & \textbf{signature (first line)} \\ \hline
— & — & — \\
\end{tabular}
\end{center}

\section*{Lean source (verbatim)}
\begin{lstlisting}
/-
  Minimal Kleisli normalization for `productMonad` (Δ ⊣ ×).
  - Stays entirely inside the Kleisli category; no mixing with TopCat arrows.
  - Uses definitional (`rfl`) lemmas so `ext; simp [Category.assoc]` cooperates
    with the η/μ component lemmas in ClaudeD.
-/

import Mathlib.CategoryTheory.Monad.Adjunction
import Mathlib.CategoryTheory.Monad.Kleisli
import MyProjects.Topology.D.ClaudeD

noncomputable section

namespace MyProjects.Topology.D

open CategoryTheory

/-- In `Kleisli productMonad`, the identity is definitional `η`. -/
@[simp] lemma kleisli_id_def (X : CategoryTheory.Kleisli productMonad) :
  (𝟙 X : X ⟶ X) = productMonad.η.app X := rfl

-- Kleisli composition is definitional: `f ≫ g = f ≫ T.map g ≫ μ`.
-- Composition in `Kleisli productMonad` is definitional:
-- `(f ≫ g : X ⟶ Z) = f ≫ productMonad.map g ≫ productMonad.μ.app Z`.
-- We omit a lemma here to avoid mixing notations; use this equality by `rfl`.

end MyProjects.Topology.D

\end{lstlisting}

\end{document}
